\uuid{jMgQ}
\titre{Matrice/Application linéaire/Noyau/Image/Rang}
\auteur{hanani ; mzari}
\organisation{exo7}
\licence{CC BY-NC-SA 4.0}
\datecreate{2020-01-20}
\niveau{L1}
\module{Algèbre}
\chapitre{Matrice}
\sousChapitre{Matrice et application linéaire}
\difficulte{1}
\theme{Matrice d'une application linéaire, Matrice, Application linéaire, Noyau, Image, Rang}
\exo7qcmid{lille-2.matrice-application.008}

\contenu{
\begin{question}
\qclassification{108.03}
\qtags{motcle=Matrice/Application linéaire/Noyau/Image/Rang}
Soit $A$ une matrice inversible d'ordre $n\ge 1$  et $f:\Rr^n \to \Rr^n$ l'application linéaire de matrice $A$ dans la base canonique de $\Rr^n$. Quelles sont les assertions vraies ?
\begin{answers}
\good{$f$ est bijective}.
\bad{Le noyau de $f$ est une droite vectorielle}.
\good{Le rang de $f$ est $n$}.
\good{Le rang de $A$ est $n$}.
\end{answers}
\begin{explanations} Les propositions suivantes sont équivalentes :
\begin{enumerate}
\item[(i)] $A$ est inversible.
\item[(ii)] Le rang de $A$ est $n$.
\item[(iii)] $f$ est bijective.
\item[(iv)] Le rang de $f$ est $n$.
\item[(v)] Le noyau de $f$ est nul.
\end{enumerate}
\end{explanations}
\end{question}
}
